\documentclass[10pt,a4paper]{amsart}
\usepackage[margin=19mm]{geometry}
\usepackage{amsmath,amssymb,mathtools}
\usepackage[scr=rsfs,cal=euler]{mathalfa}
\usepackage{xfrac}
\usepackage[unicode,hidelinks,bookmarks=false]{hyperref}
\newcommand{\ie}{i.e.\ }
\newcommand{\cf}{cf.\ }
\newcommand{\resp}{respectively\ }
\newcommand{\Subsection}{\subsection}
\providecommand{\nobreakdash}{\nobreak\hbox{-}}
\setlength{\emergencystretch}{2em}
\multlinegap0pt
\usepackage{amssymb}
\newtheorem{theorem}{Theorem}
\newtheorem{claim}{Claim}
\title{Sum of the $k$ Largest Eigenvalues of Symmetric Matrices: Theory and Applications}
\author{Shaowei Sun, Yaping Min, Kinkar Chandra Das}
\date{Source: arXiv:2605.26707v2 (CC BY 4.0)}
\begin{document}
\maketitle

\noindent\emph{Source and licence.} Sum of the $k$ Largest Eigenvalues of Symmetric Matrices: Theory and Applications. Version arXiv:2605.26707v2 (CC BY 4.0), distributed by arXiv under the Creative Commons Attribution 4.0 International licence, http://creativecommons.org/licenses/by/4.0/, as recorded in the arXiv licence metadata for this exact version and shown on the abstract page https://arxiv.org/abs/2605.26707v2. Full licence notice: \texttt{licenses/arxiv-2605.26707-CC-BY-4.0.txt}.\par\smallskip

\noindent\textit{Editorial scope.} Counted unit: the article's theorem th4 (source0.tex lines 295--305). The article's complete original proof of it is delivered with it (source0.tex lines 307--457, one \texttt{proof} environment of 151 lines). No mathematics is written, repaired or re-derived: the statement and the proof are the article's own lines, in the article's own notation, and no retained line is substituted or re-flowed.  Retained uncounted as the source-local context the proof needs: lines 174--180; lines 191--193; lines 201--203; lines 215--217; lines 232--240.  Nothing was omitted from any retained span, so the package carries no omission record.  No retained line contains a citation, so the wrapper carries no bibliography and no bibliography entry.  No retained line contains a figure environment, an image-inclusion command or drawing code, so no image is delivered and the package declares no asset dependency.  Editorial formatting only, in addition: an AMS \texttt{amsart} preamble that restates the article's own declarations and macros used by the retained lines, namely the environments \texttt{theorem}, \texttt{claim} on separate counters, together with the standard packages those lines need.  The retained environments print in this document as Theorem 1, Claim 1 in order of appearance, whereas the article prints the same items with the numbering of its own full document.  The complete native source of the article as distributed in the arXiv e-print for this exact version is bundled unchanged as \texttt{sources/201437/source0.tex}.
 \begin{theorem} \label{th1} Let $M$ be an $n\times n$ symmetric matrix. Let $s$ be an integer such that $\lambda_s(M)>\frac{tr(M)}{n}$. Then for any $k\geq s$, we have
 \begin{equation}
 \sum_{i=1}^{k}\lambda_i(M)\leq\frac{k}{n}\,tr(M)+\sqrt{tr\left(\Big(M-\frac{tr(M)}{n}I\Big)^2\right)\left(k-\frac{s^2}{\theta+s}\right)}. \label{e1}
 \end{equation}
 The equality holds if and only if  $k=s=\nu$ with 
 $$\lambda_1(M)=\lambda_2(M)=\cdots=\lambda_s(M)>\frac{tr(M)}{n}\geq \lambda_{s+1}(M)~\mbox{ and }~\lambda_{n-\theta+1}(M)=\lambda_{n-\theta+2}(M)=\cdots=\lambda_n(M).$$
 \end{theorem}

 \begin{equation}
 \sum_{i=1}^{k}\lambda_i(P)\leq \sqrt{tr(P^2)\left(k-\frac{s^2}{\theta+s}\right)}~~~\mbox{for } k\geq s, \label{e2}
 \end{equation}

\begin{align}
\sum\limits^{s}_{i=1}\lambda_i^2(P)\geq \frac{1}{s}\left(\sum\limits^{s}_{i=1}\lambda_{i}(P)\right)^2\label{1kin1}
\end{align}

 \begin{equation}
 \sum\limits^{\theta}_{i=1}\lambda_{n-i+1}^2(P)\geq \frac{\left(\sum\limits^{\theta}_{i=1}\,|\lambda_{n-i+1}(P)|\right)^2}{\theta}=\frac{\left(\sum\limits^{\nu}_{i=1}\lambda_{i}(P)\right)^2}{\theta}\geq \frac{\left(\sum\limits^{s}_{i=1}\lambda_{i}(P)\right)^2}{\theta}\label{e4}
 \end{equation}

\noindent
Otherwise, $k\geq s+1$. First we can assume that $s+1\leq k\leq \nu$. Next we will prove the inequality in (\ref{e2}) holds strictly for $s+1\leq k\leq
 \nu$. Since $\nu>s$, by Cauchy-Schwarz inequality with (\ref{1kin1}) and (\ref{e4}), we obtain
 \begin{eqnarray}
 \sum_{i=1}^{k}\lambda_i(P)&\leq&\sum\limits^{s}_{i=1}\lambda_{i}(P)+\sqrt{(k-s)\sum\limits^k_{i=s+1}\lambda_i^2(P)}\nonumber\\[3mm]
 &\leq&\sum\limits^{s}_{i=1}\lambda_{i}(P)+\sqrt{\left(k-s\right)\left(tr(P^2)-\sum\limits^{\theta}_{i=1}\lambda_{n-i+1}^2(P)-\sum\limits^{s}_{i=1}\lambda^2_{i}(P)\right)}\nonumber\\[3mm]
 &<&\sum\limits^{s}_{i=1}\lambda_{i}(P)+\sqrt{\left(k-s\right)\left(tr(P^2)-\frac{1}{\theta}\left(\sum\limits^{s}_{i=1}\lambda_{i}(P)\right)^2-\frac{1}{s}\left(\sum\limits^{s}_{i=1}\lambda_{i}(P)\right)^2\right)} \nonumber\\[3mm]
  &=&\sum\limits^{s}_{i=1}\lambda_{i}(P)+\sqrt{\left(k-s\right)\left(tr(P^2)-\frac{\theta+s}{\theta\,s}\left(\sum\limits^{s}_{i=1}\lambda_{i}(P)\right)^2\right)}. \label{e6}
 \end{eqnarray}

 \begin{theorem} \label{th4} Let $M$ be an $n\times n$ symmetric matrix. Let $s$ be an integer such that 
 $$\lambda_s(M)>\frac{tr(M)}{n}~\mbox{ and }~\sum\limits_{i=1}^s\frac{\lambda_i(M)}{s}\geq \frac{tr(M^2)+tr(M)}{n}-\frac{(tr(M))^2}{n^2}.$$ 
Then for any $k\geq s$, we have
\begin{equation}
 \sum_{i=1}^{k}\lambda_i(M)\leq\frac{k}{n}\,tr(M)+\frac{n\theta}{2(\theta+s)}+\frac{n}{2(\theta+s)}\sqrt{\theta\big(k\theta+ks-s^2\big)}.  \label{nne1}
\end{equation}
The equality holds if and only if $k=s=1$ with 
$$\lambda_1(M)=\frac{tr(M^2)+tr(M)}{n}-\frac{(tr(M))^2}{n^2}>\frac{tr(M)}{n}\geq \lambda_{2}(M)$$ 
and 
$$\lambda_{n-\theta+1}(M)=\lambda_{n-\theta+2}(M)=\cdots=\lambda_n(M).$$
\end{theorem}

\begin{proof}
If $\lambda_1(M)=\frac{tr(M)}{n}$, then $\lambda_2(M)=\lambda_n(M)=\frac{tr(M)}{n}$, which means the $s$ does not exist.
Otherwise, $\lambda_1(M)>\frac{tr(M)}{n}$, that is, $M\neq\frac{tr(M)}{n}I$ and $s\geq 1$. Let $P=M-\frac{tr(M)}{n}I$. Then we set $\nu=\nu(M)=\nu(P)$ and $\theta=\theta(M)=\theta(P)$.
 Since $tr(P^2)=tr(M^2)-\frac{(tr(M))^2}{n}$, it is equivalent to prove that for $k\geq s$,
 \begin{equation}
 \sum_{i=1}^{k}\lambda_i(P)\leq\frac{n\theta}{2(\theta+s)}+\frac{n}{2(\theta+s)}\sqrt{\theta\big(k\,\theta+k\,s-s^2\big)}~~~~\mbox{for  }\sum\limits_{i=1}^s\frac{\lambda_i(P)}{s}\geq \frac{tr(P^2)}{n},\label{ne2}
 \end{equation}
 with equality holding if and only if $k=s=1$ with $\lambda_1(P)=\frac{tr(P^2)}{n}>0\geq \lambda_2(P)$ and
 $\lambda_{n-\theta+1}(P)=\lambda_{n-\theta+2}(P)=\cdots=\lambda_n(P)$.

\vspace*{3mm}
 
 Combining Theorem \ref{th1} on $k=s$ with the condition that $\sum\limits_{i=1}^s\frac{1}{s}\lambda_i(P)\geq \frac{tr(P^2)}{n}$, we get
 $$\sum\limits_{i=1}^s\lambda_i(P)\leq\frac{n\theta}{\theta+s}\leq \frac{n\theta}{2(\theta+s)}+\frac{n}{2(\theta+s)}\sqrt{\theta\big(k\,\theta+k\,s-s^2\big)}.$$
 The first equality holds if and only if $s=\nu$, $\lambda_1(P)=\lambda_2(P)=\cdots=\lambda_s(P)=\frac{tr(P^2)}{n}$ and
 $\lambda_{n-\theta+1}(P)=\lambda_{n-\theta+2}(P)=\cdots=\lambda_n(P)$ and the second equality holds only for $k=s=1$.
 Hence we prove the case $k=s$ and together with equality cases.
 Next we will prove the inequality in (\ref{ne2}) strictly holds for $s+1\leq k\leq \nu$ as $\lambda_i(P)\leq 0$ for $i>\nu$. For this, we consider the following
two case: \vspace*{3mm}

\noindent
${\bf Case\,1.}$ $tr(P^2)\geq n\,\theta\,\sqrt{\displaystyle{\frac{tr(P^2)}{(\theta+s)(k\,\theta+k\,s-s^2)}}}=\displaystyle{\frac{n\,x_1}{s}}$ $(x_1$ is defined in the proof of Theorem \ref{th1}$)$. 

\vspace*{3mm}

In this case, $\sum\limits_{i=1}^s\lambda_i(P)\geq \frac{s}{n}tr(P^2)\geq x_1$. Recall that the function $f(x)$ defined in the proof of Theorem \ref{th1}, $f(x)$ is decreasing on $x\geq x_1$. Then from (\ref{e6}), we get
 \begin{eqnarray}
  \sum_{i=1}^{k}\lambda_i(P)&< &f\left(\sum\limits_{i=1}^s\lambda_i(P)\right)\leq f\left(\frac{s}{n}tr(P^2)\right)\nonumber\\[2mm]
  &=&\frac{s}{n}tr(P^2)+\sqrt{(k-s)\left(tr(P^2)-\frac{(tr(P^2))^2(\theta+s)\,s}{n^2\theta}\right)}\,.\label{eq1}
 \end{eqnarray}

 \vspace*{3mm}

 \noindent
 We now consider a function
 \begin{eqnarray}
 g(x)&=&\frac{s\,x}{n}+\sqrt{(k-s)\left(x-\frac{s\,(\theta+s)x^2}{n^2\theta}\right)}\,,~~~~0\leq x\leq \frac{n^2\theta}{s\,(\theta+s)}.\nonumber
 \end{eqnarray}
 Note that $g(x)$ can be rewriten as
 \begin{equation}
g(x)=\frac{s\,x}{n}+\sqrt{\frac{(s-k)(\theta+s)s}{\theta}\left(\frac{x}{n}-\frac{\theta\,n}{2s\,(\theta+s)}\right)^2+\frac{(k-s)\
     n^2\theta}{4s(\theta+s)}}. \label{sm1}
 \end{equation}
  Then we have
 \begin{equation}
 g'(x)= \frac{s}{n}+\frac{\sqrt{k-s}}{2}\times\frac{\displaystyle{1-\frac{2s(\theta+s)}{n^2\theta}x}}{\sqrt{\displaystyle{x-\frac{s(\theta+s)x^2}{n^2\theta}}}}\,,~~~0<x<\frac{n^2\theta}{s\,(\theta+s)}.\nonumber
 \end{equation}
 From $g'(x)=0$, one can easily get
 \begin{eqnarray}
  &&x_3=\frac{n^2\theta}{2s(\theta+s)}-\frac{n^2\theta}{2s(\theta+s)}\sqrt{\frac{\theta\,s}{k\theta+ks-s^2}}\nonumber\\[2mm]
  \mbox{and }&&x_4=\frac{n^2\theta}{2s(\theta+s)}+\frac{n^2\theta}{2s(\theta+s)}\sqrt{\frac{\theta\,s}{k\theta+ks-s^2}}\,.\nonumber
 \end{eqnarray}

 \noindent
 For convenience, let $x_5=\displaystyle{\frac{n^2\theta}{s(\theta+s)}}$. By the closed interval method on the continuous function $g(x)$ for $x\in \left[0,x_5\right]$, we have
 $$g(x)\leq \max\Big\{g(0),\,g(x_3),\,g(x_4),\,g_1(x_5)\Big\}.$$
 After simple calculation, we get $g(0)=0$ and $g(x_5)=\frac{n\theta}{\theta+s}$. By (\ref{sm1}), one can easily obtain that

 $$g(x_3)<g(x_4)=\frac{n\theta}{2(\theta+s)}+\frac{n}{2(\theta+s)}
 \sqrt{\frac{\theta(k\theta+ks-s^2)}{s}}.$$


 \noindent
 As $k>s$, then $\frac{k\theta+ks-s^{2}}{s}>\theta$, which follows $g(x_5)<g(x_4)$. Hence $g(x)\leq g(x_4)$.
 Applying the above results in (\ref{eq1}), we get
  \begin{eqnarray}
 &&\sum_{i=1}^{k}\lambda_i(P)< g(tr(P^2))\leq g(x_4)\leq\frac{n\theta}{2(\theta+s)}+\frac{n}{2(\theta+s)}
   \sqrt{\theta\big(k\,\theta+k\,s-s^2\big)},\nonumber
 \end{eqnarray}
 which implies that the inequality in (\ref{ne2}) strictly holds.
\vspace*{3mm}

\noindent
${\bf Case\,2.}$ $tr(P^2)<n\,\theta\,\sqrt{\displaystyle{\frac{tr(P^2)}{(\theta+s)(k\,\theta+k\,s-s^2)}}}=\displaystyle{\frac{n\,x_1}{s}}$. 

\vspace*{3mm}

As $x_1= \theta\,s\sqrt{\displaystyle{\frac{tr(P^2)}{(\theta+s)(k\,\theta+k\,s-s^2)}}}$, we have
    \begin{eqnarray}
  k\times tr(P^2)&<&\frac{n^2{\theta}^2k}{(\theta+s)(k\,\theta+k\,s-s^2)}\nonumber\\[2mm]
   &=&\frac{n^2{\theta}^2}{(\theta+s)(\theta+s-s^2/k)}\nonumber\\[2mm]
   &\leq &\frac{n^2{\theta}^2}{(\theta+s)(\theta+s/(s+1))} \label{dn1}
 \end{eqnarray}
as $k\geq s+1$. 

\vspace*{3mm}

Now let
 $$h(x)=\left(\frac{n\,x}{2(x+s)}+\frac{n}{2(x+s)}\sqrt{x\big(s\,x+x+s\big)}\right)^2-\frac{n^2x(2x-1)}{2(x+s)(x+s/(s+1))}.$$

 \vspace*{3mm}

\begin{claim}\label{c1} $h(x)>0$ for $x>0$.
\end{claim}

\vspace*{3mm}

\noindent
{\bf Proof of Claim \ref{c1}.} Now we can rewrite $h(x)$ as
   $$h(x)=\frac{n^2}{4(x+s)^2(sx+x+s)}p(x),$$ 
where
   $$ p(x)=(s^2-s-2)x^3-(2s^2-s-2)x^2+(3s^2+2s)x+2\sqrt{((s+1)x^2+sx)^3}.$$
Since $\Big((s+1)x^2+sx\Big)^3>(s+1)^2x^6$, we get 
 $$p(x)>(s^2+s)x^3-(2s^2-s-2)x^2+(3s^2+2s)x.$$
Due to the fact that 
$$(2s^2-s-2)^2-4(s^2+s)(3s^2+2s)=-8s^4 - 24s^3 - 15s^2 + 4s + 4<0$$ 
for $s\geq 1$, we conclude $(s^2+s)x^2-(2s^2-s-2)x+(3s^2+2s)$ has no real roots. By combining the above results, we obtain that $p(x)>0$ for $s\geq 1$ and $x>0$. 
Hence $h(x)>0$ for $x>0$,  which completes the proof of  ${\bf Claim\,\ref{c1}}$.
 
\vspace*{3mm}

From ${\bf Claim\,\ref{c1}}$ with (\ref{dn1}), we obtain
\begin{align}
\frac{n\,\theta}{2(\theta+s)}+\frac{n}{2(\theta+s)}\sqrt{\theta\,\Big(s\,\theta+\theta+s\Big)}&>\sqrt{\frac{n^2\theta(2\theta-1)\,(s+1)}{2(\theta+s)\Big(\theta\,(s+1)+s\Big)}}\nonumber\\[3mm]
&\geq \sqrt{k\,\left(1-\frac{1}{2\theta}\right)\,tr(P^2)}.\label{kcd1}
\end{align}

\vspace*{3mm}

Since
   $$(tr(P))^2=\left(\sum_{i=1}^{n}\lambda_i(P)\right)^2=tr(P^2)+2\sum_{i<j}\lambda_i(P)\lambda_j(P)~\mbox{ and }~tr(P)=0,$$ 
we obtain
   $tr(P^2)=-2\sum_{i<j}\lambda_i(P)\lambda_j(P)$.
Then 
    \begin{eqnarray}
 \left(\sum_{i=1}^{n}|\lambda_i(P)|\right)^2&=& tr(P^2)+2\sum_{i<j}|\lambda_i(P)\lambda_j(P)|\nonumber\\[2mm]
 &\geq &  tr(P^2)+\left|2\sum_{i<j}\lambda_i(P)\lambda_j(P)\right|\nonumber\\[2mm]
 &=&  2\,tr(P^2).\nonumber
 \end{eqnarray}
 \noindent
 Using this result, we get
  \begin{equation}
\sum_{i=1}^{\theta}\lambda_{n-i+1}^2(P)\geq
\frac{\left(\sum\limits_{i=1}^{\theta}|\lambda_{n-i+1}(P)|\right)^2}{\theta}=
\frac{\left(\sum\limits_{i=1}^{n}|\lambda_{i}(P)|\right)^2}{4\,\theta}\geq
\frac{tr(P^2)}{2\,\theta}\,.\label{e7}
 \end{equation}

 \vspace*{3mm}

 \noindent
 Using the above result with (\ref{kcd1}), we have
 \begin{eqnarray}
 \sum_{i=1}^{k}\lambda_i(P)&\leq& \sqrt{k\sum_{i=1}^{k}\lambda_i^2(P)}\nonumber\\[3mm]
     &\leq& \sqrt{k\left(tr(P^2)-\sum_{i=1}^{\theta}\lambda_{n-i+1}^2(P)\right)} \nonumber\\[3mm]
     &\leq& \sqrt{k\left(1-\frac{1}{2\,\theta}\right)\,tr(P^2)}\nonumber\\[3mm]
     &<&\frac{n\,\theta}{2(\theta+s)}+\frac{n}{2(\theta+s)}\sqrt{\theta\big(s\,\theta+\theta+s\big)}\nonumber\\[3mm]
     &\leq & \frac{n\,\theta}{2(\theta+s)}+\frac{n}{2(\theta+s)}\sqrt{\theta(k\,\theta+k\,s-s^2)}\nonumber
 \end{eqnarray}
as $k\geq s+1$. This completes the proof.
\end{proof}

\end{document}
